% SPDX-License-Identifier: CC-BY-SA-4.0
% Author: Matthieu Perrin

\begingroup

\begin{frame}[fragile]{Démontrer qu'un problème est indécidable}

  \begin{block}{Preuve par réduction}
    Supposons qu'on a déjà démontré que $P_0$ est impossible. 

    Peut-on résoudre le problème $P$ ?
    \begin{enumerate}
    \item Supposons (par l'absurde) que l'on peut résoudre $P$
    \item On prouve qu'alors on peut utiliser la solution pour résoudre $P_0$
    \item Or on sait qu'on ne peut pas résoudre $P_0$ (voir ci-dessus)
    \item Donc l'hypothèse est fausse (on ne peut pas résoudre $P$)
    \end{enumerate}
  \end{block}

  \pause

  \begin{block}{Preuve directe}
    Il faut toujours montrer qu'un premier problème $P_0$ est indécidable
    \begin{description}[labelwidth=\widthof{\bfseries Calcul}]
    \item[Géométrie plane :] la quadrature du cercle
    \item[Calcul séquentiel :] l'arrêt de la machine de Turing
    \item[Calcul réparti :] \alert{le consensus}
    \end{description}
  \end{block}

\end{frame}

\endgroup
\endinput
